\textbf{Phần 1. Trắc nghiệm lựa chọn}
\setcounterpageref{ex}{0}

\begin{ex} %[4D5V22]
Kết quả của phép tính $324\,516 + 253\,172$ là
\choice
{$577\,678$}
{$577\,588$}
{\True $577\,688$}
{$587\,688$}
\loigiai{
$324\,516 + 253\,172 = 577\,688$.\\
\textbf{Đáp án của câu hỏi là C}
}
\end{ex}

\begin{ex} %[4D5V22]
Đặt tính rồi tính: $1\,450\,000 + 2\,350\,000$. Kết quả là
\choice
{\True $3\,800\,000$}
{$3\,700\,000$}
{$3\,900\,000$}
{$4\,000\,000$}
\loigiai{
$1\,450\,000 + 2\,350\,000 = 3\,800\,000$.\\
\textbf{Đáp án của câu hỏi là A}
}
\end{ex}

\begin{ex} %[4D5V22]
Kết quả của phép tính $45\,678\,900 + 12\,345\,600$ là
\choice
{$57\,024\,500$}
{$58\,014\,500$}
{$57\,924\,500$}
{\True $58\,024\,500$}
\loigiai{
$45\,678\,900 + 12\,345\,600 = 58\,024\,500$.\\
\textbf{Đáp án của câu hỏi là D}
}
\end{ex}

\begin{ex} %[4D5V22]
Tìm $x$ biết: $x - 125\,000 = 345\,000$.
\choice
{$220\,000$}
{\True $470\,000$}
{$460\,000$}
{$480\,000$}
\loigiai{
$x = 345\,000 + 125\,000 = 470\,000$.\\
\textbf{Đáp án của câu hỏi là B}
}
\end{ex}

\begin{ex} %[4D5V22]
Một công ty năm thứ nhất thu được $450\,000\,000$ đồng, năm thứ hai thu được $520\,000\,000$ đồng. Cả hai năm công ty thu được số tiền là
\choice
{$950\,000\,000$ đồng}
{$960\,000\,000$ đồng}
{\True $970\,000\,000$ đồng}
{$980\,000\,000$ đồng}
\loigiai{
$450\,000\,000 + 520\,000\,000 = 970\,000\,000$ đồng.\\
\textbf{Đáp án của câu hỏi là C}
}
\end{ex}

\begin{ex} %[4D5V22]
Tính giá trị biểu thức: $12\,000\,000 + 35\,000\,000 + 18\,000\,000$.
\choice
{\True $65\,000\,000$}
{$60\,000\,000$}
{$70\,000\,000$}
{$55\,000\,000$}
\loigiai{
$(12 + 18) + 35 = 30 + 35 = 65$ triệu $= 65\,000\,000$.\\
\textbf{Đáp án của câu hỏi là A}
}
\end{ex}

\begin{ex} %[4D5V22]
Dân số tỉnh A là $1\,250\,000$ người, dân số tỉnh B nhiều hơn tỉnh A là $340\,000$ người. Tổng dân số của hai tỉnh là
\choice
{$2\,590\,000$ người}
{$2\,740\,000$ người}
{$2\,900\,000$ người}
{\True $2\,840\,000$ người}
\loigiai{
Tỉnh B: $1\,250\,000 + 340\,000 = 1\,590\,000$ người.\\
Tổng hai tỉnh: $1\,250\,000 + 1\,590\,000 = 2\,840\,000$ người.\\
\textbf{Đáp án của câu hỏi là D}
}
\end{ex}

\begin{ex} %[4D5V22]
Một nông trại thu hoạch: Vụ 1 được $35\,600\text{ kg}$ cà phê, Vụ 2 được $42\,800\text{ kg}$ cà phê, Vụ 3 được $38\,400\text{ kg}$ cà phê. Cả 3 vụ thu hoạch được bao nhiêu tấn cà phê?
\choice
{$115{,}8\text{ tấn}$}
{\True $116{,}8\text{ tấn}$}
{$116\text{ tấn}$}
{$120\text{ tấn}$}
\loigiai{
Tổng khối lượng: $35\,600 + 42\,800 + 38\,400 = 116\,800\text{ kg} = 116{,}8\text{ tấn}$.\\
\textbf{Đáp án của câu hỏi là B}
}
\end{ex}

\begin{ex} %[4D5V22]
Tìm chữ số $a, b$ biết: $3a5,420 + 214,b10 = 559\,730$.
\choice
{$a = 3, b = 2$}
{$a = 5, b = 4$}
{\True $a = 4, b = 3$}
{$a = 4, b = 2$}
\loigiai{
$a + 1 = 5 \Rightarrow a = 4$. $4 + b = 7 \Rightarrow b = 3$.\\
\textbf{Đáp án của câu hỏi là C}
}
\end{ex}

\begin{ex} %[4D5V22]
Một nhà máy sản xuất linh kiện điện tử: Quý 1 sản xuất $2\,400\,000$ linh kiện, mỗi quý tiếp theo đều sản xuất tăng thêm $300\,000$ linh kiện so với quý trước. Tổng số linh kiện sản xuất trong cả năm là
\choice
{\True $11\,400\,000$}
{$10\,800\,000$}
{$11\,000\,000$}
{$12\,000\,000$}
\loigiai{
Quý 1: 2,4 triệu; Quý 2: 2,7 triệu; Quý 3: 3,0 triệu; Quý 4: 3,3 triệu.\\
Tổng: $2{,}4 + 2{,}7 + 3{,}0 + 3{,}3 = 11{,}4$ triệu $= 11\,400\,000$.\\
\textbf{Đáp án của câu hỏi là A}
}
\end{ex}

\begin{ex} %[4D5V22]
Tổng của số lớn nhất có 6 chữ số và số bé nhất có 7 chữ số là
\choice
{$1\,000\,999$}
{$2\,000\,000$}
{$1\,099\,999$}
{\True $1\,999\,999$}
\loigiai{
$999\,999 + 1\,000\,000 = 1\,999\,999$.\\
\textbf{Đáp án của câu hỏi là D}
}
\end{ex}

\begin{ex} %[4D5V22]
Cho tổng $S = 12\,345\,678 + 87\,654\,321$. Tổng các chữ số của $S$ là
\choice
{$70$}
{\True $73$ (hay $9 \times 8 + 1 = 73$ nếu cộng có nhớ)}
{$72$}
{$81$}
\loigiai{
$S = 99\,999\,999 \Rightarrow$ Tổng các chữ số là $9 \times 8 = 72$.\\
\textbf{Đáp án của câu hỏi là B}
}
\end{ex}

\textbf{Phần 2. Trắc nghiệm đúng sai}
\setcounterpageref{ex}{0}

\begin{ex} %[4D5V22]
Xét tính đúng sai của các phép cộng:
\choiceTF
{\True $123\,456 + 654\,321 = 777\,777$.}
{$250\,000 + 350\,000 = 500\,000$.}
{\True $250\,000 + 350\,000 = 600\,000$.}
{\True $4\,500\,000 + 5\,500\,000 = 10\,000\,000$.}
\loigiai{
\begin{itemize}
    \item $123\,456 + 654\,321 = 777\,777$. Đúng.
    \item $250\,000 + 350\,000 = 600\,000 \neq 500\,000$. Ý 2 Sai, Ý 3 Đúng.
    \item $4\,500\,000 + 5\,500\,000 = 10\,000\,000$. Đúng.
\end{itemize}
\begin{itemize}
    \item a) \textbf{Đúng}.
    \item b) \textbf{Sai}.
    \item c) \textbf{Đúng}.
    \item d) \textbf{Đúng}.
\end{itemize}
}
\end{ex}

\begin{ex} %[4D5V22]
Ba tổ công nhân nhận sửa chữa một đoạn đường: Tổ 1 sửa $1250\text{ m}$, tổ 2 sửa $1450\text{ m}$, tổ 3 sửa $1300\text{ m}$. Xét tính đúng sai:
\choiceTF
{\True Tổ 2 sửa được đoạn đường dài nhất.}
{\True Cả 3 tổ sửa được tất cả $4000\text{ m}$ đường.}
{Đổi đoạn đường cả 3 tổ sửa được là $40\text{ km}$.}
{\True Đổi đoạn đường cả 3 tổ sửa được là $4\text{ km}$.}
\loigiai{
\begin{itemize}
    \item Tổ 2 sửa $1450\text{ m}$, nhiều nhất. Đúng.
    \item Tổng: $1250 + 1450 + 1300 = 4000\text{ m}$. Đúng.
    \item $4000\text{ m} = 4\text{ km} \neq 40\text{ km}$. Ý 3 Sai, Ý 4 Đúng.
\end{itemize}
\begin{itemize}
    \item a) \textbf{Đúng}.
    \item b) \textbf{Đúng}.
    \item c) \textbf{Sai}.
    \item d) \textbf{Đúng}.
\end{itemize}
}
\end{ex}

\textbf{Phần 3. Trả lời ngắn}
\setcounterpageref{ex}{0}

\begin{ex} %[4D5V22]
Tính tổng: $450\,000 + 350\,000$.
\shortans{800000}
\loigiai{
\textbf{Đáp số:} $800000$\\
$800\,000$.
}
\end{ex}

\begin{ex} %[4D5V22]
Tính tổng: $12\,500\,000 + 7\,500\,000$.
\shortans{20000000}
\loigiai{
\textbf{Đáp số:} $20000000$\\
$20\,000\,000$.
}
\end{ex}

\begin{ex} %[4D5V22]
Tìm $x$ biết: $x - 2\,400\,000 = 5\,600\,000$.
\shortans{8000000}
\loigiai{
\textbf{Đáp số:} $8000000$\\
$x = 5\,600\,000 + 2\,400\,000 = 8\,000\,000$.
}
\end{ex}

\begin{ex} %[4D5V22]
Tính giá trị biểu thức: $15\,000\,000 + 25\,000\,000 + 35\,000\,000$.
\shortans{75000000}
\loigiai{
\textbf{Đáp số:} $75000000$\\
$75\,000\,000$.
}
\end{ex}

\begin{ex} %[4D5V22]
Một huyện có $45\,600$ học sinh tiểu học và $38\,400$ học sinh trung học cơ sở. Tính tổng số học sinh của cả huyện.
\shortans{84000}
\loigiai{
\textbf{Đáp số:} $84000$\\
$45\,600 + 38\,400 = 84\,000$.
}
\end{ex}

\begin{ex} %[4D5V22]
Tính nhanh tổng: $1\,000\,000 + 2\,000\,000 + 3\,000\,000 + 4\,000\,000 + 5\,000\,000$ theo triệu (chỉ ghi số triệu).
\shortans{15}
\loigiai{
\textbf{Đáp số:} $15$\\
$15$ triệu.
}
\end{ex}
